\documentclass[12pt]{article}

% Packages
\usepackage[margin=1in]{geometry}
\usepackage{setspace}
\usepackage{bbm}
\usepackage{amsmath}
\usepackage{amssymb}
\usepackage{graphicx}
\usepackage{xcolor}
\usepackage{booktabs}
\usepackage{array}
\usepackage{tabularx}
\usepackage{hyperref}
\usepackage{apacite}
\usepackage{natbib}

% Paragraph spacing: blank line between paragraphs, no first-line indent
\setlength{\parskip}{0.8em plus 0.2em minus 0.2em}
\setlength{\parindent}{0pt}

% Title information
\title{Optimal Flood Policy}
\author{Taco Prins, Lennard Schlattmann, and Daniel Jonas Schmidt}
\date{\today}

\begin{document}

\maketitle

% DANIEL: this file is British spelling, draft.tex and the companion paper are American — propose going American on the merge (proper nouns stay: National Flood Insurance Program etc.).

\section{Motivation}
\label{sec:motivation}

%Flooding is the most common natural disaster in the United States, and the costliest in annual damages [reference]. Damages from flooding are projected to increase due to climate change and population movements towards exposed areas ([reference];\cite{wing2022inequitable}). 

Flooding is the most common and costliest natural hazard faced by U.S. homeowners, with losses to homeowners with federally backed mortgages alone averaging about $\$$9.4 billion per year under current climate conditions \citep{cbo2023flooddamage}. A significant increase in flood risk is projected under a wide range of climate scenarios, and adaptation by homeowners to this risk has been studied in economics as a key variable mediating the resulting increase in damages \citep{desmet2018evaluating}. 

Historically, almost all protection from flood risk enjoyed by homeowners has been provided by federal government programmes. These seek to reduce the total amount of risk, while also transferring costs from homeowners to taxpayers. Direct government assistance to homeowners in floodprone areas takes two main forms: a public flood risk insurance offer known as the National Flood Insurance Programme (NFIP) and ex-post disaster aid provided to affected households by the Federal Emergency Management Agency (FEMA). Through a mixture of incentives and regulation, the federal government furthermore encourages the building of more resilient housing in the most floodprone areas.

A key motivation of such policy interventions is that outcomes in an unfettered private market would leave homeowners exposed to excessive risks. The Act of Congress establishing the NFIP in 1968 noted that `many factors have made it uneconomic for the private insurance industry alone to make flood insurance available to those in need of such protection on reasonable terms and conditions' \citep{nfia1968}. Another problem, however, is that U.S. homeowners do not always fully take flood risk into account. This fact is very widely documented and is evident both in survey evidence \citep[e.g.][]{bakkensen2022going,mulder2024mismeasuring} and in market outcomes regarding house prices, flood insurance uptake, and flood adaptation rates. 

For example, the price discount of homes exposed to flood risk is found to be less than the discounted value of expected damages in many empirical studies\footnote{\cite{fairweather2024expecting} cite a list of 20 such studies.}; homebuyers with climate sceptic beliefs pay more for houses exposed to future flood risk through sea level rise than other buyers \citep{bernstein2019disaster,baldauf2020does,bakkensen2025leveraging}; and the discount for flood-exposed housing waxes and wanes over time as current events influence the salience of the risk in   homebuyers' minds \citep{keys2020neglected,giglio2021rfs}. \cite{wagner2022adaptation} and \cite{amornsiripanitch2025measuring} note that the uptake of flood risk insurance in high-risk Special Flood Hazard Areas (SFHAs) is relatively low (not more than 60$\%$ in each of their samples), even though insurance in these areas is mandated for many homeowners and for many properties premiums have historically been less than expected annual losses.\footnote{Furthermore, \cite{gallagher2014learning} uncovers a dynamic of Bayesian updating with forgetting that explains time-varying flood risk insurance demand at the county level as a function of nearby floods, similar to the observed salience effect of floods on house prices.} Finally, the recent empirical study of \cite{ostriker2022effects} finds that houses adapted to flooding in Florida do not command a detectable price premium, and publicly available data released by FEMA\footnote{LINK} reveal that effective adaptation measures are rare outside of areas where they are legally required.




In this paper, we formalise the trade-off faced by a policymaker dividing a fixed budget between a premium subsidy for flood risk insurance and ex-post disaster aid. We derive simple formulas for the marginal welfare gain from an extra dollar spent on either policy instrument. [1-2 further sentences describing more work.]

% DANIEL: "uninsured households are poorer" isn't in the baseline — wealth is homogeneous (w=1), so the whole marginal-utility gap is the damage channel; income is the unimplemented extension.
Some terms in our welfare formulas evaluate the transfers to different households inherent in these policies: holding insurance demand constant,  premium subsidies redistribute money to insured households, and disaster aid to uninsured households affected by floods. Households receiving disaster aid on average have higher marginal utility of consumption because floods destroy wealth and because uninsured households are poorer.

The remaining terms arise from the behavioural response to the policy instruments. First,  more generous disaster aid reduces households' use of subsidised insurance and vice versa. There is a fiscal spillover between disaster aid and premium subsidies as part of the cost of one policy is recouped through reduced spending on the other. Second, when households that underestimate flood risk are indifferent between buying insurance or not, a policy change that induces them to buy makes them better off.\footnote{The household's utility then increases by the difference between its subjective welfare estimate, conditional on being uninsured, and an objective estimate calculated using the true probability of flooding.}  This welfare gain appears in our formulas as a so-called behavioural internality. To quantify it, we posit a distribution of household beliefs over flood risk. Unlike the fiscal spillover, which depends only on the total change in insurance demand, the size of the internality follows from the location in the belief distribution where the change in insurance demand occurs.

We focus our analysis on Special Flood Hazard Areas (SFHAs), which by definition suffer a flood with an annual probability of 1$\%$ or more. The fraction of homeowners inside SFHAs with some flood insurance is much higher than in other flood-prone areas (59.5$\%$ versus 22$\%$, see \cite{amornsiripanitch2025measuring}), which reflects homeowners' greater risk awareness as well as the regulatory requirement that holders of federally backed mortgages must be insured. Our premise that homeowners maximise utility given their flood risk perception seems most appropriate to SFHAs, where it is plausible to claim that homeowners consciously decide to buy insurance or not.  

% DANIEL: this leads with the budget-share result, which implies global evaluation with fitted Beta distribution. The robust headline is the local MVPF ranking (1.33 vs 1.07, no beliefs, no lambda).
% DANIEL: "Errors in flood risk perception thus decrease the sensitivity ... to both" - I think you mean belief dispersion here, not average underestimation
Our key finding is that, over the range of plausible calibrations, optimal spending on disaster aid as a proportion of the total flood policy budget is higher than its current level. Although cumulative funding allocated to the NFIP programme has far exceeded the amount of disaster aid granted by FEMA, a substantial fraction of homeowners do not perceive flooding to be a sufficiently significant risk to purchase insurance even at a subsidised rate. This leaves disaster aid as the major safety net for this group. Moreover, the moral hazard effect of increased disaster aid is weaker when flood risk beliefs are subjective. The intuition is that households far away from the point where willingness to pay for insurance equals the actual premium are not responsive to marginal changes to the cost or benefit of insurance. Errors in flood risk perception thus decrease the sensitivity of insurance demand to both premium subsidies and disaster aid.    

% DANIEL: "runs counter to reforms that eliminate subsidies" - with the current unfinished calibration, we find s*=0 in four of six specs (mvpf_computations.md §8).
% DANIEL: We still have to settle on disaster aid parameter
In our baseline model, households vary only in their flood risk beliefs and buy insurance if their subjective flood arrival rate exceeds a threshold value. We assign values to components of the formula regarding households preferences and beliefs, the process for flood risk, and the behavioural response to policy changes to match estimates in the empirical literature. The optimal proportion of uninsured flood damages compensated by disaster aid then exceeds its current level of around 9$\%$\footnote{This number is taken from \cite{gruber2026optimal}, using their IV estimate of $\$215.70$ in disaster aid spending per flood-affected house, dividing by the application rate for disaster aid of $8\%$, and dividing the total again by the estimated average damage per house of $\$29,267$. Their definition of disaster aid encompasses not only the IHP, SBA, and HMGP programmes but also defaults on GSE-backed mortgages.} to reach between 34$\%$ and 62$\%$, depending on the assumed marginal cost of public funds ($\lambda=1.2$ and $\lambda=1.1$ respectively). The optimal premium subsidy is zero at both values of $\lambda$.\footnote{These figures come from the current, unfinished calibration described in Section~\ref{sec:calibration} and are point estimates, not bounds over the surviving calibration set; see Section~\ref{sec:results}.} This result runs counter to recent reforms to the NFIP programme that mostly eliminate premium subsidies, but is smaller than the average subsidy that pertained for most of the programme's history.

% DANIEL: Are we sure that we will implement all three extensions?
In subsequent extensions, we investigate the robustness of our central result. First, we examine the redistributive effect of expanded disaster aid across the income distribution. Second, we consider its effect on the demand for adapted housing. Third, we introduce variation in risk aversion as a further source of household heterogeneity. The qualitative result that optimal disaster aid as a fraction of flood damages exceeds its current level holds up in each extension. We also note that accounting for differences in administrative costs strengthens the case for disaster aid as a larger part of the policy mix.

% DANIEL: potential additions - charity-hazard literature, insurance demand and beliefs, my related paper
Our study is most closely related to the recent contribution by \cite{gruber2026optimal}, which analyses optimal premium subsidies holding fixed the amount of disaster aid. There, the key formula which is solved by the optimal premium subsidy is composed of empirically observable elasticities. The central estimate of the optimal flood risk insurance premium subsidy turns out to be 52$\%$, not far from the average subsidy of around 50$\%$ which according to the authors applied before the recent RR2.0 reform of the NFIP programme. Part of what renders the subsidy worthwhile is its role in dampening the reclassification risk that comes from risk-based insurance pricing. Even conditional on the amount of global atmospheric warming, projected flood risk varies greatly between climate models, and, with actuarially fair pricing, so do insurance premiums. Premium subsidies absorb the demeaned property-level component of this risk which homeowners otherwise cannot insure against.

% DANIEL: Lennard investigated complementarity of the two instruments and found the opposite of G&S, if I understood correctly
% LENNARD: Yes, I looked at the cross-partial $\partial^2 S/\partial s\,\partial a$, which was negative for a reasonable part of the belief distribution. And the joint-use gain over the better single instrument was negligible, which indicates that aid and subsidies are rather weak substitutes, opposite of G\&S. Maybe it is not only the belief addition but also the fact that they use a much higher disaster aid (13% relative to around 5% what we currently have). Probably good to recheck once we have agreed on the calibration.
The quantitative results of \cite{gruber2026optimal} are sensitive to two issues. First, in an extension where they allow for underestimation of flood risk, the optimal subsidy rises to 89$\%$, implying that the government almost fully absorbs flood risk faced by homeowners. Second, in their framework, premium subsidies and disaster aid are strongly complementary as eliminating FEMA disaster aid lowers the optimal subsidy from 52$\%$ to 15$\%$. As they write, this `joint change is welfare-improving only if households are also making rational risk and FEMA-generosity decisions' (ibid p.35), which they admit is unlikely to be the case. Importantly, we solve for  optimal public spending on premium subsidies and disaster aid without assuming either.

This paper proceeds in five parts. Section~\ref{sec:institutional} briefly describes the institutional setting determining current flood insurance provision and disaster aid spending in the United States. Section~\ref{sec:model} introduces our model of household decision-making and welfare and states the key formulas solved by the optimal levels of the premium subsidy and disaster aid. Section~\ref{sec:calibration} describes our calibration of the model to match estimates from the empirical literature. Section~\ref{sec:results} presents our results on the optimal policy mix, together with robustness checks and extensions to the baseline model. Section~\ref{sec:conclusion} concludes and offers suggestions for further research.

% As annual damages are projected to increase due to climate change, as well as population movements to vulnerable areas [references], the financial sustainability of government programmes has received the attention of an expanding economic literature.   


\section{Institutional setting}
\label{sec:institutional}

\subsection{Insurance}\label{subsec:insurance}
% DANIEL: "The NFIP programme ran a deficit almost immediately upon its establishment in 1968" - Are we sure about that? Michel-Kerjan says excluding 2005, premiums exceeded claims by ~$11bn over 1969-2008
The NFIP programme ran a deficit almost immediately upon its establishment in 1968 as on average premiums were less than expected flood damages plus administrative costs. The impact of Hurricane Katrina in 2005 effectively rendered the programme insolvent \citep{michel2010catastrophe}. The programme has also been marked by cross-subsidies from low-risk properties to high-risk properties, as premiums were set on the basis of coarse flood risk maps that did not reflect flood risk at the property level \citep{kousky2018financing}. Two recent reforms, the Biggert-Waters Act of 2012 and the Risk Rating 2.0 (RR2.0) reform of 2021-2022, have sought to eliminate insurance premium subsidies. Especially the RR2.0 reform has resulted in steeply rising premiums in parts of the United States, with rate rises spread out over several years because the maximum annual increase is capped at 18$\%$. While around 34$\%$ of policies were at the actuarially fair premium when RR2.0 was introduced, around 30$\%$ were set on a path towards the doubling of premiums or more \citep[p.~25]{gao2023flood}.  

% DANIEL: the mandate is broader than GSE (federally backed mortgages generally)
% DANIEL: We still need to settle on population for calibration and related insurance coverage number
Homeowners with a federally backed mortgage inside Special Flood Hazard Areas (SFHAs) are obliged to hold a flood insurance policy. The mandate is imperfectly enforced, however. \cite{gruber2026optimal} estimate that around 18$\%$ of all homeowners are insured when weighted by flood damages, which includes both homeowners inside and outside SFHAs. Insurance take-up is higher in SFHAs at around 50-60$\%$, but the majority of average annual losses are suffered by homeowners outside these areas \citep{amornsiripanitch2025measuring}.

% DANIEL: "actuarially fair premium" - only ~34% of policies are at full-risk rates, the rest on an 18%-capped glidepath (GAO: ~2037 before 95% get there)
Although the NFIP provides homeowners with the opportunity to insure at an actuarially fair premium, it has coverage limits of $\$$250,000 for property structures and $\$$100,000 for personal contents. These limits are occasionally binding, as shown by \cite{sastry2026bears}.\footnote{The 90th percentile of claims between 2008-2018 in the states of Florida, Texas, Alabama, Mississippi, and Louisiana  levels off at the $\$$250,000 limit for homes with an estimated replacement cost exceeding $\$$400,000, which form a non-trivial part of the distribution (ibid. Appendix Figure 1.2).} Private insurance policies, which are generally taken out in combination with an NFIP policy by owners of higher-value properties as supplementary insurance, accounted for roughly 3.5$\%$ to 4.5$\%$ of all primary residential flood policies in 2018 \citep{kousky2018emerging}.

\subsection{Disaster aid}
% DANIEL: SBA loans aren't a form of FEMA's Individual Assistance. Also "Individuals and Households Program" (plural).
When the president issues a Presidential Disaster Declaration (PDD)\footnote{Formally, the president may declare either an emergency or a major disaster under the Stafford Act, with the former releasing less funds. The SBA Administrator may also separately issue a declaration that makes SBA disaster loans available.} after a flood, individuals may apply for two main forms of so-called Individual Assistance.\footnote{Following a disaster declaration, FEMA also oversees direct transfers to states and local governments within designated counties. The Public Assistance (PA) programme covers emergency measures and repairs, while the Hazard Mitigation Grant Program (HMGP) finances projects that reduce the costs from future disasters.} First, FEMA's Individuals and Household Program (IHP) consists of grants which cover repair expenses and other immediate housing costs (Housing Assistance) and non-housing costs (Other Needs Assistance). Intended to meet `basic needs and support recovery efforts' \citep{crs2024disaster}, the maximum size of each component was $\$$44,800 as of 2026, while the average IHP grant between 2002-2024 was $\$$3,446 \citep{crs2025ihpfaq}. Second, the Small Business Administration (SBA), not directly administered by FEMA, provides loans at subsidised rates\footnote{The annual interest rate of SBA loans may not exceed 4$\%$ `if the applicant is found 
by SBA to be unable to obtain credit elsewhere' \citep[p.~2]{crs2024disaster}.} with maturities of up to 30 years. These loans are capped at $\$$500,000 for home repairs and $\$$100,000 for personal property replacement. 

Homeowners with federally backed mortgages form a special category of aid recipients. A disaster declaration that releases Individual Assistance funding also makes homeowners eligible for forbearance for up to 12 months on their mortgage repayments for GSE-backed loans.\footnote{Information accessed via this \hyperlink{https://www.fhfa.gov/homeowners-and-homebuyers/mortgage-assistance/disaster-assistance}{link}.} Furthermore, it has been argued that mortgage lenders' ability to sell securitised mortgages to GSEs encourages them not to adjust loan conditions to flood-related default risk, giving homeowners access to underpriced implicit insurance against flood damages  \citep{keenan2020underwaterwriting,sastry2023insurers}. \cite{bakkensen2025leveraging} and \cite{liaomulder2026homeequity} find evidence that homeowners exploit this insurance option in a strategic manner.\footnote{\cite{bakkensen2025leveraging} show that buyers identified as pessimistic regarding climate change are more likely to use a mortgage when purchasing properties exposed to sea level rise and have higher loan-to-value ratios,. \cite{liaomulder2026homeequity} show that owners of highly leveraged properties are less likely to purchase flood risk insurance.} However, \cite{gallagher2017household} note that Hurricane Katrina engendered at most a modest rise in defaults, while \cite{sastry2026bears} demonstrates that mortgage lenders restrict credit in flood-prone areas to avoid excess losses, mainly by reducing maximum loan-to-value ratios.\footnote{\cite{nguyen2022climate} provide earlier evidence of reduced LTVs in flood-prone areas, but, using data , note that providers underestimate.} 


\subsection{Adaptation}

Federal legislation encourages and in some cases mandates flood risk adaptation. Besides compensating for behavioural biases and misaligned incentives within the housing market, the federal government thereby addresses to some extent the moral hazard that its own provision of subsidised insurance and disaster aid creates. 

Communities participating in the NFIP programme must enforce a minimum set of regulations within floodplains. The most important of these is the requirement, in force since the roll-out of the NFIP in the 1970s and 1980s, that homes constructed within SFHAs be elevated beyond the projected height of a 100-year flood (the Base Flood Elevation). This measure reduces the average annual loss from flooding by around 56$\%$ (\cite{ostriker2022effects}). It is relatively inexpensive at the time of construction, although costs for existing properties are generally forbidding (\cite{botzen2019adoption}).   

The NFIP's Community Rating System (CRS) provides points-based discounts on insurance premiums to communities that implement further adaptation measures ranging from public information provision to drainage improvement. Weighted by the number of policies per community, participation in the CRS is high at around 70$\%$. However, it is also shallow as most participating communities earn modest discounts through low-cost measures (\cite{fema2021crs}).

Several federal programmes, such as FEMA's Hazard Mitigation Program, allow communities to access subsidies to finance flood protection structures. Public flood protection measures are less extensive in the United States than in other flood-prone developed countries (e.g., the Netherlands) and only about 1.5$\%$ of homeowners live in levee-protected areas (\citealp{qin2026flood}).

\section{The model}
\label{sec:model}

We formalise the trade-off described above in a one-period model of a single flood-prone region. A social planner allocates a fixed budget $B$ between two policy instruments: an insurance subsidy $s$ as a fraction of the actuarially fair flood insurance premium, and disaster aid $a$ as a fraction of uninsured flood damages. Households differ only in their subjective beliefs about flood risk, which affects their binary decision whether or not to buy insurance.

\subsection{Environment}

A continuum one of households each own a coastal property and hold wealth $w$. They face a flood event that occurs with true, common probability $p$. Expressed as a proportion of wealth, the damage conditional on a flood event  follows a known distribution with mean $\overline{d}$ and CDF $G(d)$.

Households hold heterogeneous subjective beliefs $q$ about the probability that a flood event occurs, distributed across the population according to the CDF $F(q)$. By contrast, the distribution of damages conditional on flooding $G(d)$ is known and accepted by all households. If a household buys insurance, it pays the insurance premium $(1-s)p\overline{d}$, since the actuarially fair premium is $p\overline{d}$. Flooding is an idiosyncratic risk affecting a continuum one of households, so the actuarially fair premium faced by the household at the start of the period exactly equals realised damages paid out by the insurer at the end of the period.

\subsection{Household problem}

We assume that household chooses between no insurance and full insurance. An insured household's consumption is therefore unaffected by whether a flood occurs. Its utility is
\begin{align}
    V_I(s) = u(c_I), \quad \text{with} \quad c_I = w - (1-s)p\overline{d}.
\end{align}
An uninsured household consumes $c_{U,N} = w$ if no flood occurs and $c_{U,F}(d) = w-(1-a)d$ if it does. The social planner evaluates the expected utility of the household using the true flood probability $p$, so
\begin{align}
    V_U(a) &= (1-p)\,u(c_{U,N}) + p\,\mathbb{E}[u(c_{U,F}(d))] = u(c_{U,N}) - p\cdot\Delta u,
\end{align}
where $\Delta u \equiv u(c_{U,N})-\mathbb{E}[u(c_{U,F}(d))]>0$ is the expected utility loss an uninsured household suffers from a flood. Meanwhile, the household with subjective belief $q$ perceives its expected utility from remaining uninsured to be
\begin{equation}
    \tilde{V}_U(q,a) = u(c_{U,N}) - q\cdot\Delta u.
\end{equation}
The household voluntarily buys insurance whenever $V_I(s) > \tilde{V}_U(q,a)$, that is, whenever its subjective belief exceeds a threshold $q^*(s,a)$ implicitly defined by indifference between the two options:
\begin{equation}
    q^*\cdot\Delta u = u(c_{U,N}) - u(c_I). \label{eq:qstar_prog}
\end{equation}
Furthermore, given the legal requirement that homeowners with a federally backed mortgage hold flood risk insurance, we assume that with probability $\phi$ the household is required to buy insurance independently of its subjective belief $q$.\footnote{In reality, compliance with the regulatory requirement falls over the lifetime of the mortgage as insurance needs to be renewed annually. No precise estimate of overall non-compliance exists, and we ignore the issue.}  

Given the CDF for flood risk beliefs $F(q)$, the insurance take-up rate is $I(s,a) = \phi+(1-\phi)(1-F(q^*(s,a)))$. The threshold $q^*(s,a)$ itself is independent of the distribution of beliefs.

\subsection{Social planner problem}

Households make decisions that maximise their subjective expected welfare, and, given those decisions, the social planner evaluates welfare using the true flood probability $p$. Social welfare is the population average of household utility,
\begin{equation}
    S(s,a) = V_I(s)\cdot I(s,a) + V_U(a)\cdot\bigl(1-I(s,a)\bigr), \label{eq:welfare_prog}
\end{equation}
which the planner seeks to maximise subject to the budget constraint that total spending on insurance subsidies and disaster aid not exceed $B$:
\begin{equation}
    B = p\overline{d}\bigl[s\,I(s,a) + a\,\bigl(1-I(s,a)\bigr)\bigr]. \label{eq:budget_prog}
\end{equation}

\subsubsection{MVPF given current policy $(s,a)$}

Attaching multiplier $\lambda$, the shadow value of the budget constraint, to~\eqref{eq:budget_prog} and differentiating the resulting Lagrangian with respect to $s$ and $a$ yields two first-order conditions. Rearranging each so that the welfare gain per dollar of fiscal cost equals $\lambda$ gives the marginal value of public funds (MVPF) of the subsidy and of disaster aid,
\begin{equation}
    \text{MVPF}_s = \frac{
        \dfrac{\partial V_I}{\partial s}\cdot I
         +  (V_I-V_U)\,\dfrac{\partial I}{\partial s}
    }{
        p\overline{d}\!\left[I + (s-a)\,\dfrac{\partial I}{\partial s}\right]
    }, \qquad
    \text{MVPF}_a = \frac{
        \dfrac{\partial V_U}{\partial a}\cdot(1-I)
         + (V_I-V_U)\,\dfrac{\partial I}{\partial a}
    }{
        p\overline{d}\!\left[(1-I) + (s-a)\,\dfrac{\partial I}{\partial a}\right]
    }. \label{eq:mvpf_prog}
\end{equation}
For the planner's choice $(s,a)$ for the policy instruments to reach an interior optimum, it must be that $\text{MVPF}_s = \text{MVPF}_a = \lambda$. If that condition is not satisfied, the planner should (marginally) increase the budget share of whichever instrument has the higher MVPF.

Each numerator combines two channels through which the corresponding instrument affects welfare. The first is the value of the transfer to households who do not change their decision as a result of the marginal change in policy: the subsidy benefits every insured household by $\partial V_I/\partial s = p\overline{d}\,u'(c_I)$ per unit, while aid benefits every uninsured household by $\partial V_U/\partial a = p\mathbb{E}[d\,u'(c_{U,F}(d))]$ per unit. Because utility is concave (marginal utility is decreasing in consumption), Jensen's inequality guarantees that $V_U/\partial a >V_U/\partial s $. Holding household decisions fixed, disaster aid thus increases social welfare by more than the insurance subsidy. 

The second channel takes account of the households which are induced by the policy change to switch between insurance and no insurance. Households hold mistaken beliefs, so the envelope theorem does not apply. Instead, households' behavioural responses to marginal policy changes affect their own welfare. Each switcher generates a behavioural internality that scales with the difference between the true probability $p$ and the belief of the marginal household $q^*$:
\begin{equation}
    V_I - V_U = \tilde{V}_U(q^*,a) - V_U(a) = (p-q^*)\,\Delta u, \label{eq:internality_prog}
\end{equation}
where we follow the internality terminology of \cite{allcott2019regressive}. The utility loss from flooding $\Delta u$ is positive, so the internality as a whole is positive provided that $p>q^*$. Given that the insurance premium before subsidies is actuarially fair and households are risk-averse, this requires that $a$ does not exceed $s$ by too much. In section \ref{sec:results}, we find that the internality is positive for current policy and for all $(s,a)$ we consider as candidates for optimal policy. Therefore, in line with the policy debate about underinsurance referenced in Section ~\ref{sec:motivation}, increased insurance take-up improves social welfare. The behavioural internality term creates a case for the insurance subsidy (if $\partial I/\partial s>0$) and against disaster aid  (if $\partial I/\partial a<0$).   

The denominator is the marginal cost of each policy instrument. This sums, firstly, the direct transfer from each instrument to the households it targets and, secondly, the cost from switching households net of the fiscal spillover between the instruments. To understand the latter, note that if an increased insurance subsidy causes a household to purchase insurance, the cost to the government is only $(s-a)p\overline{d}$ because the household will no longer require disaster aid once it suffers a flood (and, similarly, expanding disaster aid saves on subsidies).    



Note that, so far, we have not made any assumptions about the difficult-to-observe distribution $F(q)$ of flood risk belief across households. The terms in \eqref{eq:mvpf_prog} are current policy $(s,a)$, a process for flooding, assumptions regarding household wealth and risk aversion, and finally empirically observable insurance demand. We are thus able to calibrate each term and obtain the marginal value of each policy instrument while remaining agnostic about the distribution of household beliefs (although clearly not about the belief $q^*$ of the marginal household). Furthermore, the parameter $\phi$ does not appear in the terms, since it does not matter to the social planner whether homeowners purchase insurance voluntarily or otherwise.

\subsubsection{MVPF given optimal policy $(s^*,a^*)$}
Away from current policy $(s,a)$, the empirical values of the behavioural terms in \eqref{eq:mvpf_prog} obtained from recent data (namely, the insurance take-up $I$ and its derivatives w.r.t. the policy instruments) are no longer valid. To find the optimal combination of policy instruments $(s^*,a^*)$, we must posit a distribution of flood risk beliefs $F(q)$. This distribution then implies the insurance take-up rate $I=\phi+(1-\phi)(1-F(q^*(s,a)))$ and its marginal response to policy.


The chain rule tells us that 
\begin{equation}\label{eq:partialsinsurance}
    \frac{\partial I}{\partial s} = -(1-\phi)f(q^*)\,\frac{\partial q^*}{\partial s}, \qquad
    \frac{\partial I}{\partial a} = -(1-\phi)f(q^*)\,\frac{\partial q^*}{\partial a},
\end{equation}
while from the indifference condition~\eqref{eq:qstar_prog} 
\begin{equation}\label{eq:partialsqstar}
    \frac{\partial q^*}{\partial s} = -\frac{p\overline{d}\,u'(c_I)}{\Delta u} < 0, \qquad
    \frac{\partial q^*}{\partial a} = \frac{q^*\,\mathbb{E}[d\,u'(c_{U,F}(d))]}{\Delta u} > 0.
\end{equation}
The response of insurance uptake to the two policy instruments is not symmetric. The reduction in the insurance premium from the subsidy $s$ scales with the true probability $p$ (which enters the actuarially fair premium) and therefore so does the marginal household's response to the subsidy. By contrast, the value placed by the marginal household on disaster aid depends on its subjective probability $q^*$. To the extent that $q^*$ lies below $p$, its response to disaster aid is attenuated. On the other hand $\mathbb{E}[d\,u'(c_{U,F}(d))]>\overline{d}u'(c_I)$\footnote{If not, disaster aid is (more than) large enough that the insurance premium exceeds the expected damage from flooding conditional on it occurring. We do not find this to be the case for current or optimal policy.}, so whether $\partial I/\partial s>\partial I/\partial a$ or vice versa is a priori ambiguous.


Substituting these responses, together with $\partial V_I/\partial s$, $\partial V_U/\partial a$, and the internality~\eqref{eq:internality_prog}, into~\eqref{eq:mvpf_prog} and cancelling the common factors $p\overline{d}$ and $\Delta u$ yields the marginal value of public funds in terms of the model's deep parameters:
\begin{align}
    \text{MVPF}_s &= \frac{u'(c_I)\bigl[I + (p-q^*)\,(1-\phi)f(q^*)\bigr]}{I + (s-a)\,\dfrac{p\overline{d}\,u'(c_I)\,(1-\phi)f(q^*)}{\Delta u}}, \notag \\
    \text{MVPF}_a &= \frac{\dfrac{\mathbb{E}[d\,u'(c_{U,F}(d))]}{\overline{d}} \! \left[(1-I) - \dfrac{(p-q^*)q^*(1-\phi)f(q^*)}{p}\right]}{(1-I) - (s-a)\,\dfrac{q^*d\,\mathbb{E}[d\,u'(c_{U,F}(d))]\,(1-\phi)f(q^*)}{\Delta u}}. \label{eq:mvpf_deep_prog}
\end{align}

These terms summarise the arguments rehearsed above. The numerator of each term describes the utility effect of the policy and the denominator its cost. Net of any behavioural effect, disaster aid has a stronger effect than the insurance subsidy on utility ($\mathbb{E}[d\,u'(c_{U,F}(d))]/\overline{d}>u'(c_I)$). Meanwhile, the behavioural effect of the subsidy on utility is positive and that of disaster aid is negative (the terms in the numerator containing the belief wedge $(p-q^*)$). However, the behavioural effect of disaster aid is dampened given that $q^*<p$ (both in the numerator and denominator of the second term). Finally, in the denominator, the fiscal spillover is represented by the difference $(s-a)$ as before. 

As mentioned, the terms in \eqref{eq:mvpf_deep_prog} can only be quantified given a belief distribution $F(q)$. Although we cannot observe this distribution directly, clearly a minimum requirement of a candidate distribution is that it ensures \eqref{eq:mvpf_prog} and \eqref{eq:mvpf_deep_prog} are consistent: given current policy $(s,a)$, the insurance take-up $I$ should match its observed level, and the density $f(q^*(s,a))$ should reproduce the marginal response of insurance demand to policy which we insert directly into \eqref{eq:mvpf_prog} using empirical estimates. We describe our calibration strategy for both equations in greater detail below. We will refer to the terms in \eqref{eq:mvpf_prog} as the `current policy MVPF' and to the terms in \eqref{eq:mvpf_deep_prog} as the `optimal policy MVPF'. 

%Each numerator is the marginal utility of the instrument's direct recipients, $u'(c_I)$ or $u'(c_{U,F})$, scaled by an effective recipient count: the inframarginal share of households already making that choice, plus a belief-gap correction contributed by switching households equal to their probability mistake $(p-q^*)$, weighted by the density $f(q^*)$ of households located at the margin. The correction is smaller for aid by a factor $q^*/p<1$, reflecting the weaker behavioural response documented above. When no households are located exactly at the margin, $f(q^*)=0$, all behavioural terms vanish, and each instrument's marginal value collapses to the marginal utility of consumption of its direct recipients: $\text{MVPF}_s=u'(c_I)$ and $\text{MVPF}_a=u'(c_{U,F})$.

%The denominators mirror this structure. Each combines the direct fiscal cost of supporting the inframarginal group, $I$ for the subsidy and $1-I$ for aid, both measured relative to $pd$, with a fiscal externality term equal to $(s-a)$ times the deep-parameter take-up response derived above: $pd\,u'(c_I)\,f(q^*)/\Delta u$ for the subsidy, and $q^*d\,u'(c_{U,F})\,f(q^*)/\Delta u$ for aid. The same asymmetry that governs the internality carries over here: the subsidy's fiscal externality scales with the true flood probability $p$ acting on the premium, while aid's scales with the marginal household's own subjective probability $q^*$ acting on its flood-state consumption. Where $q^*<p$, as in our calibration, aid's fiscal spillover onto the subsidy budget is muted relative to the mirror-image spillover the subsidy exerts on the aid budget, for the same reason its take-up response is muted: households near the bottom of the belief distribution are simply less responsive to policy at the margin.

%Two empirical objects are therefore sufficient to evaluate both MVPFs at the current policy: the insurance take-up rate $I$, and the density $f(q^*)$ of households at the insurance margin. We recover the latter from the empirically estimable elasticity of take-up with respect to the insurance premium, without needing to observe the belief distribution $F$ directly. Evaluating counterfactual policies away from the status quo, however, requires the full distribution $F(q)$; we obtain this by fitting a parametric distribution to match the observed take-up rate and premium elasticity. Section~\ref{sec:calibration} describes this calibration in detail.

\subsection{The dynamic problem}

In the one-period problem developed thus far, damages from flooding are wholly subtracted from current consumption. When there are multiple periods, households are able to maintain their consumption by borrowing or drawing down their assets. In the United States, government policy seeks to assist homeowners in . Homeowners with a GSE-backed mortgage enjoy an automatic forbearance (not forgiveness) of mortgage payments for a period of 6 to 12 months in areas receiving FEMA assistance. This temporarily loosens households' borrowing constraint after a flooding disaster, in order to prevent mortgage defaults and blunt the immediate drop in consumption. Here, we examine how consumption smoothing alters the cost of uninsured flooding disasters and thereby our formulas for optimal policy.   

\subsubsection{Household problem}

There is an infinite number of periods, and the flooding disaster with damage distribution $G(d)$ now strikes each period with constant probability $p$. As before, households are coastal homeowners and their belief about the flood arrival rate $p$ is described by the distribution function $F(q)$.  They now choose anew each period whether to buy flood insurance at cost $\overline{P}=(1-s)p\overline{d}$. If they are hit by a flood while uninsured, they suffer damages $(1-a)d$.

Households' per period-utility $u(c)$ is a function of consumption $c$ (not wealth $w$). They receive stochastic income $e$ each period which follows a Markov chain, have identical discount factor $\beta$, and can borrow and save at interest rate $r$.  



In a reduced-form way, we incorporate in our model the limited ability of households to access credit immediately following a flooding disaster. Let $B_{\text{max,U,F}}(w,e)$ be the maximum amount an uninsured household is able to borrow after it suffers a flood; $B_{\text{max,U,I}}(w,e)$ the same if it is insured or it does not suffer a flood. It is also helpful to define the cash-in-hand that results from the insurance decision and the flood realisation:

\begin{equation}
\begin{aligned}
m_\text{I} &= w + e - \overline{P}, \\
m_\text{U,F}   &= w + e - (1-a)d, \\
m_\text{U,N}  &= w + e.
\end{aligned}
\end{equation}

Finally, let $c_{\text{max}}$ be the resulting maximum consumption level, e.g., $c_{\text{max,I}}(w,e)=m_\text{I}(w,e)+B_{\text{max,I}}(w,e)$. 


The household with subjective belief $q$ regarding the flood arrival rate maximises the subjective expectation of its welfare. The timing of each period is as follows: the household chooses whether to insure, the flooding process is realised, and finally the household chooses consumption. Given its start-of-period wealth $w$, income level $e$, its problem is given by 

\begin{equation}
\tilde{V}(w,e;q)=\max\left(\tilde{V}_{I}(w,e;q),\,\tilde{V}_{U}(w,e;q)\right),
\end{equation}
where
\begin{equation}
\tilde{V}_{\text{I}}(w,e;q)=\max_{c\leq c_{\text{max,I}}}\; u(c)+\beta \mathbb{E}_{e'}[\tilde{V}((1+r)(m_\text{I}-c),e';q)]
\end{equation}
and
\begin{align}
\tilde{V}_{\text{U}}(w,e;q)= q\mathbb{E}_d[\max_{c\leq c_{\text{max,U,F}}}\; u(c)+\beta \mathbb{E}_{e'}[\tilde{V}((1+r)(m_\text{U,F}-c),e';q)]\\
+(1-q)\max_{c\leq c_{\text{max,U,N}}}\; u(c)+\beta \mathbb{E}_{e'}[\tilde{V}((1+r)(m_\text{U,N}-c),e';q)]
\end{align}


%\begin{equation}
%\tilde{V}(w,e;q)=\max\left(\tilde{V}_{ins}(w,e;q),\,\tilde{V}_{unins}(w,e;q)\right),
%\end{equation}
%where
%\begin{equation}
%\tilde{V}_{ins}(w,e;q)=\max_{c\leq c_{\text{max,ins}}}\; u(c)+\beta \mathbb{E}_{e'}[V((1+r)(w+e-c-\overline{P}),e';q)]
%\end{equation}
%and
%\begin{align}
%\tilde{V}_{unins}(w,e;q)= q\mathbb{E}_d[\max_{c\leq c_{\text{max,f}}}\; u(c)+\beta \mathbb{E}_{e'}[V((1+r)(w+e-c-(1-a)d)),e';q)]\\
%+(1-q)\max_{c\leq c_{\text{max,nf}}}\; u(c)+\beta %\mathbb{E}_{e'}[V((1+r)(w+e-c),e';q)]
%\end{align}

Let $i$ be the resulting policy function for the insurance decision and $c$ the policy function for consumption conditional on insurance and the realisation of flooding. Now, define 
\begin{equation}
\Delta b(d,w,e;q)=1-\frac{c_\text{N}(d,w,e;q)-c_\text{F}(w,e;q)}{(1-a)d}.
\end{equation}

$\Delta b$ is the additional amount the household borrows or withdraws from savings in response to a flood event, as a fraction of the damage it suffers. If the household is already close to its borrowing constraint ($c_\text{U,N}\approx c_\text{max,U,N})$ and the amount of additional credit available to a household after a flood is low ($B_\text{max,U,F}\approx B_\text{max,U,N}$), $\Delta b$ is close to $0$. Current consumption is then reduced by the flood damage $(1-a)d$. Conversely, if the borrowing constraint is not binding, $\Delta b$ is much closer to $1$ and the household does not suffer a consumption trough.      

The policy functions together with the processes for income and flooding define a stochastic process for consumption $c_t$. The social planner calculates the objective value of the household from this flow: 
\begin{equation}
V(w,e;q)=\mathbb{E}\!\left[\textstyle\sum_{t\geq0}\beta^{t}\,u(c_t)\;\middle|\;w_0=w,\,e_0=e\right],
\label{eq:objective_dyn}
\end{equation}
where the expectation is taken over the income process and the true flood arrival rate $p$.




%The objective (true-probability) value of a household, given the decisions generated by its
%subjective problem, is obtained by evaluating the same policy functions under the true arrival
%rate $p$:
%\begin{equation}
%V(w,e;q)=\mathbb{E}\!\left[\textstyle\sum_{t\geq0}\beta^{t}\,u(c_t)\;\middle|\;w_0=w,\,e_0=e\right],
%\label{eq:objective_dyn}
%\end{equation}
%where the expectation is taken over the true income process and the true flood arrival rate $p$,
%but consumption $c_t$ and the insurance choice $\mathbbm{1}_{\text{ins}}$ follow the subjective
%policy functions defined above. The wedge between $V$ and $\tilde{V}$---decisions taken under the
%belief $q$, welfare scored under $p$---is the source of the behavioural internality, exactly as in
%the static model of Section~\ref{sec:model}.

\subsubsection{The approximation technique}

In the one-period model, the cost of flooding was given by $\Delta u \equiv u(c_{U,N})-\mathbb{E}_d[u(c_{U,F}(d))]$ with $c_{U,F}(d)=c_{U,N}-(1-a)d$. In the dynamic model, the cost depends on the household's ability to smooth its consumption. The household evaluates the expected cost of a flood as \begin{align}
\Delta \tilde{V} \equiv u(c_{U,N})-\mathbb{E}[u(c_{\text{U,N}}-(1-\Delta b)(1-a)d)]+ \nonumber \\ \beta\mathbb{E}_{e'}([\tilde{V}((1+r)(m_\text{U,N}-c_\text{U,N})]-\mathbb{E}_{d}[\tilde{V}((1+r)(m_\text{U,N}-c_\text{U,N}-\Delta b(1-a)d)]),
\end{align}
where the dependence of objects on the states $(w,e)$ and belief $q$ is suppressed for readability. The social planner uses the same formula to evaluate the cost of flooding to the household, albeit with the objective value $V$ replacing $\tilde{V}$.

We now introduce a simple idea: since flood damages are generally small compared to lifetime earnings (but not one-period consumption), then, provided the household is not at its borrowing constraint next period, next-period value functions $\tilde{V}$ and $V$ are approximately linear over the range given by flood damages $(1-a)d$. In this case, we can write:
\begin{align}\label{eq:approxformula}
\Delta \tilde{V} \approx u(c_{U,N})-\mathbb{E}_d[u(c_{U,N}-(1-\Delta b)(1-a)d)]+ \nonumber \\ \beta(1-a)(1+r)\mathbb{E}_d[\Delta b(d)d]\cdot\mathbb{E}_{e'}([\tilde{V}_w((1+r)(m_\text{U,N}-c_\text{U,N})]).
\end{align}

Now consider the case where the uninsured household is close to its borrowing constraint before the flood. In that case, the value function clearly has curvature when future wealth changes by $\Delta b(1-a)d$, but we expect  $\Delta b$ to be close to $0$. Therefore, equation \eqref{eq:approxformula} continues to hold approximately. 

We can simplify the formulas further by using $u'(c_{U,N})=\beta(1+r)\mathbb{E}_{e'}([\tilde{V}_w((1+r)(m_\text{N}-c_\text{N})$ (again, this works on the second term, which is relevant if the household is not credit-constrained). Then,   

\begin{equation}\label{eq:approxformulasimple}
\Delta \tilde{V} \approx u(c_{U,N})-\mathbb{E}[u(c_{U,N}-(1-\Delta b)(1-a)d)]+\beta(1-a)\mathbb{E}_d[\Delta b(d)d]\cdot u'(c_{U,N}),
\end{equation}

and, since the insurance premium is small,

\begin{equation}
\tilde{V}_{\text{N}}-\tilde{V}_{\text{I}}\approx\overline{P} \cdot u'(c_{U,N}).
\end{equation}

The indifference condition to buy insurance is now given by

\begin{equation}
q^*\approx\frac{\overline{P}}{(u(c_{U,N})-\mathbb{E}_d[u(c_{U,N}-(1-\Delta b)(1-a)d)])/u'(c_{U,N})+(1-a)\mathbb{E}_d[\Delta b(d)d]}.
\end{equation}

To interpret this condition, note that the denominator contains two parts. The first part concerns the hit to current consumption and contains a risk premium, since $u(c_{\text{U,N}})-u(c_{\text{U,F}})>u'(c_{\text{U,N}})(1-\Delta b)(1-a)d$ if the agent is risk-averse. The second part concerns the part of flood damages that the agent is able to spread out over periods. In the extreme case where the agent can smooth consumption perfectly ($\Delta b\approx1$), the agent becomes approximately risk-neutral and $q^*$ equals the insurance premium divided by expected flood damages. This follows from the fact that flood damages are small relative to lifetime earnings. 

\subsubsection{Implications for the MVPF formulas}

The structure of the social planner problem described in \eqref{eq:welfare_prog} and \eqref{eq:budget_prog} is unchanged, except that the social planner eventually has to aggregate over agents by their wealth and income level. Setting this problem aside for now,  

%In a reduced-form way, we incorporate in our model the limited ability of households to access credit immediately following a flooding disaster. Let $\mathbbm{1}_{\{\text{insured}\}}$ and  $\mathbbm{1}_{\{\text{flood}\}}$ be indicators describing whether the household has flood insurance whether a flood has occurred. Then, let $b_{\text{max}}(w,e,\mathbbm{1}_{\{\text{insured}\}},\mathbbm{1}_{\{\text{flood}\}})$ be the maximum amount the household is able to borrow after the realisation of the flood process. $c_\text{max}(w,e,\mathbbm{1}_{\{\text{flood}\}})=w+e+b_{\text{max}}(w,e,\mathbbm{1}_{\{\text{insured}\}},\mathbbm{1}_{\{\text{flood}\}})-(1-a)d$ is the amount the household can consume each period. 

%In a reduced-form way, we incorporate in our model the limited ability of households to access credit immediately following a flooding disaster. Let $d_{\text{net}}$ be the flood damage the household suffers after the realisation of the flood process. It is zero if the household is insured or if it does not suffer a flood; otherwise it is $d_{\text{net}}=(1-a)d$. Then, let $b_{\text{max}}(w,e,d_{\text{net}})$ be the maximum amount the household is able to borrow after the realisation of the flood process and $c_\text{max}(w,e,d_{\text{net}})$ be the maximum amount the household can consume in the period. If an insured household suffers a flood, $c_\text{max}(w,e,d_{\text{net}})=w+e+b_{\text{max}}(w,e,d_{\text{net}})-(1-a)d$.


%The simple assumption underpinning our reduced-form approach is that, when hit by a flood, households are not able to divide the loss in wealth over current and future periods in an unconstrained fashion. Instead, in a period in which a household suffers a flood, a maximum fraction of flood damages $1-b(d)$ may be borrowed. The remainder $b(d)$ must be drawn from savings or is subtracted from current consumption. The reduced-form object $b(\cdot)$ represents the sum of unmodelled forces affecting households' ability to smooth their consumption: on the one hand, loan-to-value and payment-to-income constraints as well as delays in obtaining credit in disaster-stricken areas, and on the other hand government policy expanding the availability of credit.      




\subsection{How to identify belief density $f(q^*)$}

From Gruber and Solomon, we know the response of insurance uptake to a percentage increase in the premium: 
\begin{equation}
\epsilon\equiv\frac{\partial I}{\partial\pi/\pi}\approx-0.32.
\end{equation}

We can integrate this semi-elasticity in our model as follows. Since the premium $\pi=(1-s)pd$, $\partial \pi/\partial s=-\pi/(1-s)$. Therefore,
\begin{equation}
\frac{\partial I}{\partial s}=-\frac{\partial I}{\partial\pi/\pi}\frac{1}{1-s}.
\end{equation}

Combining with \eqref{eq:partialsinsurance} and \eqref{eq:partialsqstar} gives

\begin{equation}
f(q^*)=\frac{1}{1-\phi}\frac{-\epsilon}{\pi}\frac{\Delta u}{u'(c_I)}
\end{equation}


\section{Calibration}
\label{sec:calibration}

\subsection{Current policy MVPF}



\subsubsection{The flooding process}


To do's
\begin{itemize}
\item Convert distribution of damages as $\%$ of house value to wealth loss as in model. How? Using simple toy model with MPV? If so should be introduced in section 3, and made consistent with welfare formulas...
\item For flood arrival rate -- given the distribution of damages, make consistent with AAL reported in \cite{amornsiripanitch2025measuring}?
\end{itemize}

\subsubsection{Utility}

\subsubsection{Current policy $(s,a)$}

In our model, $s$ is the wedge between the flood insurance premium paid by the household and its expected insurance pay-out. It is not (minus) the balance of the NFIP programme divided by its revenues, which includes administrative costs that do not affect the household. Administrative costs are an important consideration in practice, and we discuss them in section \ref{subsec:admincosts}.

As mentioned in section \ref{subsec:insurance}, current policy is not stable, because the RR2.0 reform has set insurance premiums on a gradual path towards actuarially fair rates. Based on the availability of high-quality empirical studies, we calibrate our `current policy MVPF' formulas from data just \textit{before} the NFIP reform; accordingly, the value of $s$ is positive.

There is a range of estimates reported in the literature for the NFIP subsidy rate, some ex-ante (comparing premiums to estimated flood risk) and some ex-post (comparing premiums to realised flood events). The ex-post estimate of \cite{wagner2022adaptation} is a subsidy of 30 per cent. The ex-ante estimate of \cite{gruber2026optimal}, which comes simply from calculating the average nationwide premium increase under RR2.0, is even larger at 46 per cent. \cite{bakkensen2020sorting} produce a similar ex-ante estimate of about 50 per cent in a sample of Florida properties, reflecting pre-FIRM and CRS discounts built into the system prior to RR2.0. We accordingly set $s=0.46$. 

For the disaster aid fraction $a$, there is a careful estimate from \cite{gruber2026optimal}. They employ an instrumental variables strategy in which exogenous variation in the county-level insurance rate predicts federal spending after a flood event. Their estimate of $a=0.133$ includes spending under the HWBP programme, which goes to local authorities and not to households. Removing this component, we arrive at $a=0.12$.\footnote{\cite{gruber2026optimal} also account for losses on federally backed mortgages following a disaster. It is not obvious how this fits our framework as the foreclosure loss to the lender may not equal the household's financial gain from default. As a component of disaster aid, it is in any event quantitatively unimportant.}  


\subsubsection{Insurance take-up}
Although almost 59.5$\%$ of homeowners within SFHAs hold some flood insurance, only 48.4$\%$ of expected flood damages are insured as not all homeowners purchase full insurance (see \cite{amornsiripanitch2025measuring}; statistics as of 2023). In our model with a binary insurance decision, the aggregate insurance rate equals the percentage of expected damages covered. We therefore set $I=0.484$, including both voluntary and mandated insurance demand. 

Regarding the fraction $\phi$, \cite{mulder2024mismeasuring} finds in his survey of homeowners in high-risk areas that 39$\%$ are subject to a mandatory insurance requirement. He uses that number to convert an estimate for the price elasticity of insurance demand to an estimate for \textit{voluntary} insurance demand. For our purposes, an important adjustment comes from the fact that the purchase requirement is partial.\footnote{It is capped at the minimum of the value of the mortgage, the NFIP coverage limit of $\$250,000$, and the replacement value of the property} Unfortunately, there is no data to compare the intensive margin of the insurance decision between homeowners with and without a mortgage.\footnote{Researchers have been unable to match individual-level mortgage data with publicly available NFIP policies and claims data because the latter are anonymised; \cite{mulder2024mismeasuring}'s survey avoids this problem.} We therefore assume their policies on average cover the same fraction of expected damages and set $\phi=0.39\cdot\frac{0.484}{0.595}=0.32$. It follows that the \textit{voluntary} insurance take-up rate inside SFHAs is approximately 24.5$\%$, which is significantly higher than the rate of 16.3$\%$ outside SFHAs.\footnote{See \cite{amornsiripanitch2025measuring}, p. 969: in flood-prone areas outside SFHAs, 22$\%$ of households hold a policy and 16.3$\%$ of expected damages are covered.}    

Several estimates exist in the literature for the price elasticity of flood insurance demand. \cite{mulder2024mismeasuring} exploits NFIP map adjustments transferring homeowners into high-risk SFHAs, which are accompanied by price increases at different calendar dates. He arrives at a price elasticity of insurance demand of $-0.177$ (and $-0.37$ for voluntary demand). \cite{wagner2022adaptation}, using price variation from the Biggert-Waters and HFIAA reforms of the NFIP, obtains a price elasticity of $-0.25$ in SFHAs. Finally, \cite{gruber2026optimal} calculate a semi-elasticity of $-0.32$ from the renewal of policies subject to price rises after the RR2.0 reform. For a baseline results, we will assume a price elasticity of $-0.21$, in the middle of the first two estimates which exploit recent and exogenous price variation within SFHAs.         

We map from these empirical elasticities to the derivatives in \eqref{eq:mvpf_prog} as follows. If we label the effective premium faced by households $\overline{P}$, then $\overline{P}=(1-s)p\overline{d}$ and $\partial \overline{P}/\partial s = -\overline{P}/(1-s)$. Therefore,
\begin{equation}\label{eq:insurancedemandwrts}
\frac{\partial I}{\partial s}\bigg|_{(s,a)}=\frac{\partial I}{\partial \overline{P}}\frac{\partial  \overline{P}}{\partial s}\bigg|_{(s,a)}=-\frac{\partial I/I}{\partial\overline{P}/\overline{P}}\frac{I}{1-s},
\end{equation}
where the first fraction on the RHS is the elasticity. Filling in our estimates for this elasticity and for $I$ and $s$ leads to $\partial I/\partial s=0.19$. It should be emphasised again that this is a local estimate that holds given current policy $(s,a)$. Meanwhile, from \eqref{eq:partialsinsurance} and \eqref{eq:partialsqstar}, 
\begin{equation}\label{eq:insurancedemandwrta}
\frac{\partial I}{\partial a}=-\frac{q^*}{p}\frac{\mathbb{E}[d\,u'(c_{U,F}(d))]}{p\overline{d}}\frac{\partial I}{\partial s}
\end{equation}


\section{Results}
\label{sec:results}

\section{Robustness}
\label{sec:robustness}
\subsection{Redistributive effects}
Households vary in their wealth and income, which bears on their exposure to flooding and their demand for flood insurance. Thus, the insurance premium subsidy and disaster aid may  have redistributive effects.

\cite{qin2026flood} combine comprehensive census and property tax for the whole of the United States to document a number of facts. First, moving across the income distribution, the fraction of households living in a floodplain is essentially flat, hovering between around 13$\%$ in the left tail to just under 12$\%$ around the 75th percentile, with however a spike in the top income decile reaching around 18$\%$ (ibid p. 31). Second, relative to income, both homeowners' insurance premiums and their expected losses from flooding fall considerably with income, and with a comparable slope (ibid p. 44 and p. 46). Third, the insurance rate rises with income -- moving from the 25th to the 75th income percentile, the fraction of uninsured households in floodplains falls by about half (ibid p. 45).\footnote{Compared to the administrative NFIP data of \cite{amornsiripanitch2025measuring}, the survey data of \cite{qin2026flood} indicate a higher insurance rate overall.} 

The effects of the premium subsidy and of disaster aid on income inequality seem minor: both instruments do little for the poorest households which do not own housing but tend to favour less wealthy homeowners, depending on the manner in which they are financed. If raw exposure to flooding does not change very much with income, the fact that the flood insurance rate rises with income would make the insurance subsidy less progressive than disaster aid. To the extent that redistributive concerns affect the marginal value of either instrument, we therefore expect our baseline result (given current policy, $\text{MVPF}_a>\text{MVPF}_s)$ to be strengthened.

\subsection{Demand for adapted housing}

\subsection{Administrative costs}\label{subsec:admincosts}

Here, enter some available estimates. 

\subsection{Functional form belief distribution}

\subsection{Dispersion in risk aversion}
\section{Conclusion}
\label{sec:conclusion}

\bibliographystyle{apacite}
\bibliography{bibfile.bib}
\end{document}